\documentclass[twoside]{article}
\usepackage[utf8]{inputenc}

% The package supports other formats as well, but this template
% is primarily checked and updated for examinations at TUM BGDM
%\usepackage[uebung, answers]{tumbgdm}

\usepackage[uebung]{tumbgdm}
\usepackage{pythonhighlight}

\Professor{Prof. Dr. Martin Werner}
\Universitaet{Technical University of Munich}
\Institut{School of Engineering and Design}
\Lehrstuhl{Professorship Big Geospatial Data Management}
\Semester[WiTe 2024/25]{Winter Term 2024/25}
\Vorlesung{Principles of Programming}
\InstitutsWebsite{https://www.bgd.lrg.tum.de/}
\VorlesungsWebsite{https://www.moodle.tum.de/course/view.php?id=103704}

\nummerUebungsblatt{4}
\nummerErsteFrage{1}

\LearningOutcome{ Problem solving}

\begin{document}
\maketitle

\vspace{10pt}
The following problems are meant to be solved with pure Python and do not require any additional packages. 
You must recall all the previously discussed concepts (lists, loops, conditions, etc.) to solve them.
However, you may optionally also try to solve them with any other programming language of your choice.

\begin{task}{Single Number}{}{}
Given a non-empty array of integers \code{nums}, every element appears twice except for one. Find that single one.
\vspace{10pt}

\textbf{Example 1:}
\begin{lstlisting}[numbers=none]
Input: nums = [2,2,1]
Output: 1
\end{lstlisting}
\vspace{10pt}

\textbf{Example 2:}
\begin{lstlisting}[numbers=none]
Input: nums = [4,1,2,1,2]
Output: 4
\end{lstlisting}
\vspace{10pt}

\textbf{Example 3:}
\begin{lstlisting}[numbers=none]
Input: nums = [1]
Output: 1
\end{lstlisting}

\end{task}

\clearpage

\begin{task}{Sum Multiples}{}{}
Given a positive integer \code{n}, find the sum of all integers in the range $[1, n]$ inclusive that are divisible by $3$, $5$, or $7$.
Return an integer denoting the sum of all numbers in the given range satisfying the constraint.
\vspace{10pt}

\textbf{Example 1:}
\begin{lstlisting}[numbers=none]
Input: n = 7
Output: 21
\end{lstlisting}
Explanation: Numbers in the range [1, 7] that are divisible by 3, 5, or 7 are 3, 5, 6, 7. The sum of these numbers is 21.
\vspace{10pt}

\textbf{Example 2:}
\begin{lstlisting}[numbers=none]
Input: n = 10
Output: 40
\end{lstlisting}
Explanation: Numbers in the range [1, 10] that are divisible by 3, 5, or 7 are 3, 5, 6, 7, 9, 10. The sum of these numbers is 40.
\vspace{10pt}

\textbf{Example 3:}

\begin{lstlisting}[numbers=none]
Input: n = 9
Output: 30
\end{lstlisting}
Explanation: Numbers in the range [1, 9] that are divisible by 3, 5, or 7 are 3, 5, 6, 7, 9. The sum of these numbers is 30.

\end{task}



\begin{task}{Count Prefixes}{}{}
You are given a string array \code{words} and a string s, where every \code{words[i]} and \code{s} comprise only lowercase English letters.
Return the number of strings in words that are a prefix of \code{s}.
A prefix of a string is a substring that occurs at the beginning of the string. 
A substring is a contiguous sequence of characters within a string.

 \vspace{10pt}

\textbf{Example 1:}

\begin{lstlisting}[numbers=none]
Input: words = ["a","b","c","ab","bc","abc"], s = "abc"
Output: 3
\end{lstlisting}
Explanation: The strings in words which are a prefix of \code{s} = "abc" are:
"a", "ab", and "abc".
Thus the number of strings in words which are a prefix of \code{s} is 3.
\vspace{10pt}

\textbf{Example 2:}
\begin{lstlisting}[numbers=none]
Input: words = ["a","a"], s = "aa"
Output: 2
\end{lstlisting}
Explanation: Both strings are a prefix of \code{s}. 
Note that the same string can occur multiple times in words, and it should be counted each time.

\end{task}

\end{document}